\section{Conclusiones y Trabajos Futuros}

\subsection{8.1. Conclusiones clave}

Resulta revelador observar cómo las diferencias entre arquitecturas de combustión, aunque sutiles en algunos parámetros, se acumulan hasta definir trayectorias distintas para la aviación comercial sostenible. El combustor anular demostró consistentemente superioridad en eficiencia de combustión, uniformidad térmica y reducción de emisiones, con mejoras de NOx cercanas al 35-40\% y mejor comportamiento en perfiles de temperatura que benefician la durabilidad de la turbina.

Jeong y Ahn ya anticipaban parte de estas ventajas en sus análisis unidimensionales de double annular combustors, donde la optimización de parámetros geométricos y de flujo conducía a reducciones significativas tanto de NOx como de CO sin penalizar excesivamente la estabilidad \cite{Jeong2021}. 

Nuestros resultados confirman y extienden esos hallazgos bajo condiciones representativas de motores de alto bypass. No obstante, los combustores can-anulares mantuvieron ventajas claras en robustez transitoria y facilidad de mantenimiento, recordándonos que el desempeño en papel no siempre se traduce directamente en viabilidad operativa diaria.

\subsection{8.2. Recomendaciones de diseño}

Uno de los desafíos persistentes que surge al cerrar este análisis radica en traducir las ventajas técnicas en decisiones concretas de ingeniería. Para motores de nueva generación se recomienda priorizar arquitecturas anulares cuando los objetivos principales sean reducción agresiva de emisiones y eficiencia en crucero, especialmente en combinación con SAF. 

La implementación debe ir acompañada de sistemas de enfriamiento avanzados y tolerancias de fabricación más estrictas. En cambio, para actualizaciones de motores existentes o aplicaciones donde el mantenimiento y la certificación rápida resulten críticas, los can-anulares siguen ofreciendo un camino más seguro y económico.

Particularmente llamativo es el hecho de que configuraciones híbridas o modulares podrían combinar lo mejor de ambos mundos: liners parcialmente segmentados dentro de una estructura predominantemente anular. Los diseñadores deberían además incorporar desde las primeras etapas consideraciones de ciclo de vida completo, incluyendo costos de producción y potencial de reparación en taller.

\subsection{8.3. Trabajos futuros}

Lejos de ser un asunto resuelto, el camino hacia la propulsión aeronáutica net-zero exige investigación continua y multidisciplinaria. Chen et al. enfatizan que el verdadero potencial de los combustibles sostenibles solo se materializará cuando las arquitecturas de combustión estén específicamente optimizadas para sus características físicas y químicas \cite{Chen2025}.

Entre las líneas prioritarias destacan: validación experimental en instalaciones de alta presión de los perfiles de emisiones con mezclas SAF superiores al 50\%; exploración de combustión de hidrógeno en geometrías anulares con énfasis en mitigación de flashback y oscilaciones termoacústicas; desarrollo de modelos híbridos CFD-1D que incorporen aprendizaje automático para optimización rápida de parámetros; y estudios económicos detallados que cuantifiquen el trade-off real entre costo inicial y beneficios operativos y ambientales a lo largo de la vida útil del motor.

Adicionalmente, sería valioso investigar configuraciones activamente controladas con inyección adaptativa y sensores en tiempo real, así como el impacto de degradación por envejecimiento en el desempeño comparativo de ambas arquitecturas. 

Este trabajo ha iluminado tendencias claras y trade-offs relevantes. Sin embargo, solo mediante colaboración estrecha entre academia, industria y autoridades regulatorias lograremos cerrar la brecha entre simulaciones prometedoras y soluciones certificadas listas para surcar los cielos de 2035 y más allá.